% Revised Abstract (ontology paragraph aligned with thesis implementation)
\chapter{Abstract}

This thesis investigates the limits and opportunities of learning robust morphologic representations from lung adenocarcinoma histopathology under realistic clinical constraints: small cohorts, class imbalance, inter-observer disagreement, and intrinsically fuzzy morphologic boundaries.

Rather than treating histologic pattern classification as an end in itself, we reformulate it as an intermediate representation learning problem whose ultimate objective is the inference of underlying genetic alterations. We systematically benchmark a family of margin-based and fuzzy margin-based loss functions, with particular focus on successive variants of FuzzyArcLoss, to identify embedding geometries that maximize stability and generalization under morphologic ambiguity.

Across extensive ablation studies on the Zenodo ANORAK lung adenocarcinoma dataset (775 samples, 6 classes), we observe a consistent performance ceiling around 76--77\% macro-F1 for patch-level histologic classification. The best-performing configuration, FuzzyArcLoss v3 with sub-centers, explicitly models intra-class multimodality and achieves 76.96\% macro-F1 and 76.07\% accuracy. Advanced strategies---including curriculum scheduling, ensemble learning, and evolutionary optimization---fail to yield consistent improvements beyond this ceiling, indicating an epistemic limitation of hard patch-level supervision rather than an optimization bottleneck.

To connect local morphology with global tumor biology, patch-level predictions are aggregated into case-level morphologic profiles capturing the relative prevalence of architectural growth patterns. These aggregated representations are used to predict clinically relevant genetic mutations (EGFR, TP53, KEAP1, STK11) using gradient-boosted decision trees, with SHAP-based explainability revealing mutation-specific morphologic drivers.

A key finding is that mutations associated with global architectural disruption, particularly TP53, are more robustly inferred from aggregated morphologic composition than mutations linked to subtle epithelial differentiation, such as EGFR. This behavior aligns with established genotype--phenotype relationships in lung adenocarcinoma and provides biological validation for the aggregation-based design of the pipeline.

Finally, an ontology-grounded explanation layer links model outputs to standardized biomedical knowledge through a three-tier design: a curated knowledge graph over canonical IRIs (NCIt, MONDO) with provenance from WHO, OncoKB/FDA, and TCGA/CIViC; dynamic case-specific subgraphs that retain only patterns, genes, and treatments relevant to each patient; and a literature-driven enrichment pipeline (DeepSearch) that discovers and validates new relations from PubMed, Semantic Scholar, and arXiv. SPARQL retrieval and an LLM-based natural-language summary provide auditable, evidence-traceable explanations without influencing training or prediction.
